% !TEX root = handbook-physics.tex
% !TEX TS-program = pdflatex
% !TEX encoding = UTF-8 Unicode






\chapter{Diffraction of Light}
\label{ch-33}

\section{Huygens-Fresnel Principle}
\label{sec-33.1} 

\subsection{}\label{33.1.1} If the position of the wavefront (30.2.3) is known at a certain
instant t of time, as well as the velocity v of the wave, the
position of the wavefront at the subsequent instant t + At can
he determined on the basis of the Huygens principle. According
to this principle, each point on the surface S (t), through which
the wavefront passes at the instant t can be regarded as a source
of secondary waves, and the sought-for position S (t + At) of
I he wavefront at the instant of time t + At coincides with the
envelope of all the secondary waves at this later time. It is
jissumed here that in a homogeneous medium the secondary waves
radiate only forward, i.e. in directions making acute angles
with the outward normals to the wavefront. In a homogeneous
isotropic medium, the secondary waves are spherical (Fig. 33.1).

\subsection{}\label{33.1.2} The Huygens principle can be used to derive laws for the
reflection and refraction of light at the interface between two
media. The line MN in Fig. 33.2 represents the plane interface
between two media in which the velocities of light equal vx
and u2. Incident on this interface at the angle i is a plane wave
(rays 1 and 2). At the instant t of time the wavefront (plane
AB) reaches the interface at point A. Hence point A begins
to radiate secondary waves, propagating both in the first medium
(reflected wave) and in the second medium (transmitted wave).
During the time At required for the incident wave to travel the distance BC (wher A/ = BC!v ), the front of the secondary
wave, radiated fr< i point A, reaches points of a semisphere in
the first medium having the radius Rf = i>xA* = BC, and points
of a semisphere in the second medium having the radius R2 =
= u2At = v BClvi, The front ol the reflected wave (rays 1'
V A t £
N
Fig. 33.1 Fig. 33.2
and 2'), propagating at the angle V of reflection is the plane DC,
tangent to a sphere of radius R± with its centre at point A.
Correspondingly, the front of the transmitted (refracted) wave
(rays 1" and 2"), propagating at the refraction angle r, is the
plane CE, tangent to a semisphere of radius R2 with its centre
at point A, The law of light reflection', i = i' follows from the
equality of triangles A CD and ACB. The law of light refraction
follows from the rectangular triangles ACB and ACE, which
have a common hypotenuse. Thus
sin i
BC
Vf
 
   
 
sin r ~ AE ~ vt ~ ”21’
where n2t is the relative refractive index of the second and
first media (31.5.1).


\subsection{}\label{33.1.3} The Huygens principle is a purely geometrical one. It
indicates no method of calculating the amplitude of the wave
that is an envelope of the secondary waves. Therefore, the
Huygens principle is insufficient for deriving the laws of propaga-tion of light waves. An approximate method of solving this
problem, which is a development of the Huygens principle on
the basis of the idea proposed by Fresnel on the coherence of
secondary waves and their interference upon superposition is
called the Huy gens-Fresnel principle. This principle can be
expressed by the following series of propositions:
(a) In calculating the amplitude of light vibration, set up by
source S0, at arbitrary point M, the source can be replaced by
a system of secondary sources. These sources are small portions
ds of any closed auxiliary surface S arranged so that it encloses
the source S0J but does not enclose point M being considered.
(b) The secondary sources are coherent with S0 and with one
another. Hence, the secondary waves they set up interfere upon
superposition. Interference calculations are simplest when S
is the wavefront for the light from source S0 because, in this
case, the phases of vibration of all secondary sources are the
same.
(c) The amplitude dA of the vibration set up at point M by
a secondary source is proportional to the ratio of the area ds
of the corresponding portion of wavefront S to the distance r
from this portion to point M and depends upon the angle a
between the outward normal to the wavefront and the direction
of a line from the element ds of area to point M.
Thus
1 A ( / \   dA~-f(a)-y- %
where a is a quantity proportional to the amplitude of the primary
wave at points of the element ds of area; / (a) decreases
monotonically from 1 at a = 0 toO at a > jx/2 (secondary
sources do not radiate backward)*.
(d) If part of surface S is an opaque screen, the corresponding
secondary sources (closed by the screen) do not radiate, whereas
the rest radiate in the same way as in the absence of screens.**
* As shown by Kirchhoff, / (a) oc (1 + cos a), i.e. the function
becomes zero only at a = jx; at small angles of diffraction a,
however, this more exact relation is of no significance.
** Actually, the material of the opaque screen has an effect
on the radiation of the open secondary sources that are close to
the edge of the screen (at a distance of the order of the wavelength
of the light).

\subsection{}\label{33.1.4} On the basis of the Huygens-Fresnel principle, the law
of rectilinear propagation of light in a homogeneous medium can
be substantiated from the wave point of view. Assume that S0 is
a point source of monochromatic light (Fig. 33.3) and M is the
point of observation. As the auxiliary surface S, a spherical
surface of radius R is chosen so that the distance L from point M
to this sphere (L = OM) is of the order of R. Next, surface S
is divided into annular portions of small area, called Fresnel
zones, as shown in Fig. 33.3, where X is the wavelength of the
light.
Vibrations set up at point M by two adjacent zones are opposite
in phase, because the path difference from like points of these
zones to point M is equal to X/2. Consequently, the amplitude of
the resultant vibration at point M is equal to A = — A 2 +
+ As — A4 + . . where At is the amplitude of vibration set
up at point M by secondary sources, located within the limits of
the ith zone. At £<c LlX, the areas of all the zones are the same:
oi = JtRLX/(R + L), and very small (at R = L = 10 cm and
X = 5 X 10"6 cm we obtain   8 X 10"5 cm2). With an
increase in i, the distance rt from the zone to point M also
increases, as does angle at between the normal to the surface of
the zone and the direction of a line from the zone to point M.

Therefore, according to the Huygens-Fresnel principle, A1 >»
> A 2 > A3 > . . . and At   ( i-i + Ai+1)/2. Consequently,
the amplitude of vibration at point M is equal to A   Aj2y
i.e. the resultant effect of the whole open wavefront is equal to
<»ne-half of the effect of the first (central) Fresnel zone, whose
radius is very small. Thus, it can be assumed that practically
light is propagated from S0 to M in a straight line.


\subsection{}\label{33.1.5} If an opaque screen, covering all, except the first, Fresnel
zones for the point M of observation is placed in the path of
monochromatic light from source S0, the amplitude and intensity
of light at point M are doubled and quadrupled, respectively,
compared to their values in the absence of the screen: A=At
and / = A\. Light can be intensified to a substantially greater
extent by means of a zone plate. This is a glass plate to which an
opaque coating has been applied in the form of rings which
enclose only the even (or only the odd) Fresnel zones, i.e. every
other zone. A zone plate has the same effect on light as a converging
lens.

\section{Fresnel Diffraction}
\label{sec-33.2} 


\subsection{}\label{33.2.1} The diffraction of light is the whole set of phenomena due
lo the wave nature of light and observed in its propagation in
a medium with sharply defined optical inhomogeneities (for
example, in passing through apertures in opaque screens or
( lose to the boundaries of opaque bodies). In its narrower sense,
diffraction refers to the bending of light around small opaque
obstacles, i.e. deviations from the laws of geometrical optics.
Distinction is made between two kinds of light diffraction:
Fresnel diffraction, or diffraction with convergent rays, and Fraunhofer
diffraction, or diffraction with parallel rays. In the former,
a spherical or plane wave is incident on the obstacle and the
diffraction pattern is observed on a screen at a finite distance
behind the obstacle. In the latter, a plane wave is incident on
I lie obstacle and the diffraction pattern is observed on a screen
located in the focal plane of a converging lens placed in the
path of the light that passes by the obstacle. In Fresnel diffrac-
t ion, a “diffraction image” of the obstacle is obtained on the
screen; in Fraunhofer diffraction, the “diffraction image” is of the
distant light source.

\subsection{}\label{33.2.2} In the simplest Fresnel diffraction problems, the diffrac-tion pattern can be investigated by using the method of Fresnel
zones (33.1.4).
Example 1. Fresnel diffraction by a small circular aperture in an
opaque screen AB (Fig. 33.4). When the aperture is illuminated
by monochromatic light of wavelength X, a system of alternating
dark and bright interference rings is observed on screen E,
which is parallel to AB. These rings have a common centre at
point Oy located opposite the centre of the aperture. If for point 0
an even number 27c (where 7c = 1, 2, . . .) of Fresnel zones fit into
the aperture, there is a dark spot at point 0\ the amplitude of
light at point 0 is less than in the absence of a screen:
A (AA2k)   A\ /2»
where At is the amplitude corresponding to only the ith zone. If
the number of zones is odd (2k - - 1), there is a bright spot
at point 0 and
A   (A1-A2h+i) > AJ2.
If the aperture is illuminated by white light, a system of concentric
coloured rings is observed on screen E.
The number of Fresnel zones that fit into the aperture and the
degree of contrast of the interference pattern depend upon the
ratio of the diameter d of the aperture to the distance I between
screens AB and E. As djl is increased, the amplitude of light atthe centre of screen E approaches A j2 and the degree of contrast
of the interference rings is reduced.

\subsection{}\label{33.2.3} Example 2. Fresnel diffraction by a small disk (opaque
round screen). The method of constructing open Fresnel zones
on the wavefront S of the incident monochromatic spherical
wave is illustrated in Fig. 33.5. The interference pattern on
screen E consists of alternating bright and dark rings with their
centre at point 0, where the interference maximum (Poisson s
spot) is always observed. The amplitude of light at point 0 is
one-half of the amplitude Alt corresponding to the effect at this
point of only the single first open Fresnel zone: A = Aj2.
When the disk is illuminated by white light, a white spot is
observed at the centre of the screen. This spot is surrounded by
a system of concentric coloured rings.
As the ratio of the diameter d of the disk to its distance I from
screen E is increased, the brightness of Poisson’s spot is gradually reduced and the next dark ring becomes broader, forming a region
of shadow beyond the disk.

\section{Fraunhofer Diffraction}
\label{sec-33.3} 


\subsection{}\label{33.3.1} Example 1. Diffraction of light by a narrow long slit
in an opaque screen (Fig. 33.6). The slit width BC — b, and its
length in the direction perpendicular
to the plane of • the figure is I > b.
Light is normally incident on the
surface of the slit so that vibrations
at all points of the slit are in phase.
The diffraction pattern is observed on
screen E, which is placed in the focal
plane of converging lens L. Parallel
rays BM and CN, travelling from the
edges of the slit at the diffraction
angle if) to the direction of incident
light, are collected by the lens at its
secondary focus F . The lens are made
so that the paths of the rays of light
BMFyp and DNF , where D is the
Fig. 33.6
foot of the perpendicular dropped
from point B to ray CN, are tautochronous paths (32.3.2).
Therefore, the result of the interference at point of the screendepends upon the path d i fference CD = b sin and the wavelength
X of the light.
(a) Approximate solution. The slit can be divided in width into
Fresnel zones, having the form of strips parallel to edge B, such
that the path difference from the edges of the strips is equal
to X/2. The number of such Fresnel zones that fit into the slit
is equal to 2 6 | s i n A l l the zones radiate light entirely
equally in the direction being considered; and the vibrations
set up at point by two adjacent zones are equal in amplitude
and opposite in phase. Therefore, when there is an even
number of zones:
6 sin of) = ± 2 m y , where ra = l, 2, ... ,
and a diffraction minimum is observed (complete darkness).
When the number of zones is odd:
X
b sin )= ± (2m +1) — , where m = l, 2, ... ,
and a diffraction maximum corresponding to one Fresnel zone,
is observed. The brightest central maximum is observed at the
principal focus F0 of the lens (ip = 0). With an increase in m the
width of the Fresnel zones and the intensity of the maxima are
rapidly reduced.
(b) Exact solution. The slit is divided into a very great number
of identical very narrow strips, parallel to edge B. Secondary
waves radiated by these elements of the slit, set up vibrations at
point Fx j, that have the same small amplitudes, and their
initial phases continuously fill an interval of width Atp =
= (2nb sin ty)/X. According to Sect. 32.4.4, the amplitude and
intensity of light at point equal:
Tib sin ob
sin --- t— — sin2----t——
. 0 Tib sin ob
where A0 and I0 are the amplitude and intensity at the central
maximum (o|) = 0). 'The condition for diffraction minima of light is the same as in
the approximate solution: b sin ? — ±2mX/2, where m =
— 1,2,... . The condition for diffraction maxima is of the
form
Tib sin  
nb sin 'll?
 
and it differs only slightly from the condition obtained by means
of the Fresnel zones method.


\subsection{}\label{33.3.2} If white light, instead of monochromatic light, falls on
the slit, the central maximum is white with iridescent colouring
along its edges. All the other interference fringes are coloured
because different angles 'll? and different points on the screen
correspond to both minima and maxima of the same orders m,
depending upon the wavelength X.
As the width b of the slit is reduced, the width of the central maximum is increased because of the increase of angles   =
= + arc sin (k/b), which correspond to minima of the first order,
limiting the central maximum. AX b .k the illumination of the
screen is reduced monotonically from the centre (point F0)
toward the edges. When the slit is very wide (b> k), a bright
and sharp image of Ithe light source, formed by lens L according
to the laws of geometrical optics, is seen on the screen.


\subsection{}\label{33.3.3} Example 2. Diffraction of light by a circular aperture.
A plane monochromatic light wave is normally incident on the
aperture, i.e. perpendicular to the plane of the aperture. The
diffraction pattern is observed in the focal plane of a converging
lens arranged beyond the aperture so that its optical axis is
perpendicular to the plane of the aperture. The diffraction pattern
has the form of a bright white spot at the principal focus F0
of the lens and, concentric about this spot, alternating dark and
bright rings. The intensity of the bright rings is very low compared
to that 10 of the central maximum, and decreases with an
increase in their radius. For instance, the intensity of the nearest
maximum of the first order to the central one is Ix < 0.02 /0.
The diffraction angle i|?l5 corresponding to the first dark ring,
bounding the central maximum, complies with the condition where D is the diameter of the aperture and X is the wavelength
of the light.
If light falls on the aperture at a small angle a to the normal to
the plane of the aperture, the interference pattern remains practically
unchanged, but its centre is displaced to a secondary focus of the lens, corresponding to the angle   = a.


\subsection{}\label{33.3.4} Example 3. Diffraction of light by a one-dimensional
grating. A one-dimensfbnal diffraction grating is a system of
iarge number N of parallel
slits in a screen, the slits being
\ tal tal ifi i
identical in width and separated
by opaque strips, also of equal
width. Shown in Fig. 33.7 are
only two adjacent slits of the
grating. The quantity d= a-\-b,
where a = CD is the width of an
opaque strip between two slits
and b = BC is the slit width, is
called the diffraction grating constant,
or interval.
In analysing the diffraction pattern
on screen E, mounted in
the focal plane of converging
lens L, it is necessary to take
into account the interference of
the secondary waves not only
from different portions of one slit, but from the different slits
as well. If a plane monochromatic wave is normally incident
on the grating, the vibrations at all points of the slits are in
phase.
Vibrations set up at the arbitrary point of the focal plane
of lens L by each of the slits coincide in amplitude (Aj) and
differ in phase. For each pair of adjacent slits the phase difference
Acp0 between these vibrations is the same. It depends upon
the wavelength X and the path difference from like points of the
slits (for instance, points B and D), i.e. on the length KD =
= d sin “ip, where K is the foot of the perpendicular, dropped
from point B to the ray DN. Thus
2nd
A<Po =
sin i|).Therefore, according to Sects. 32.4.2 and 33.3.1, the amplitude
and intensity of the resultant vibration at point F  equal
sin itb sin ?
sin nNd sin i|?
%
%
A=Af
nb sin *i|?
nd sin 'll?
Kb sin a)? sm  KNd sin  
r = r n
( Kb sin "ij?  2
I I )Kd sin oj?
where A0 and 70 are the amplitude and intensity of the vibration
at point F0 (i.e. at i|? = 0) due to the effect of a single slit.
\subsection{}\label{33.3.5} The principal minima in the diffraction of light by
a diffraction grating occur at diffraction angles i|?, corresponding
!o interference minima in diffraction by a single slit (33.3.1).
Thus
h sin 'll? = + (m = 1, 2, . . .).
In these directions jione of the slits provide light (“each slit
extinguishes itself”).
The p r in cip a l m a x im a (32.4.3) correspond to diffraction angles'll?
that comply with the condition
d sin i|? = + nk,
where n = 0, 1, 2, ... is the order of the p r in cip a l m a x im u m .
Note. If certain values of a|? simultaneously comply with the
conditions for both principal maxima and principal minima,
then the principal maxima corresponding to these values of o|?
are not observed. For example, if d = 2b, then all the even principal
maxima (n = 2, 4, 6, etc.) are absent.
The intensity of a principal maximum of the rath order is
T ( Nd \ 2 . 2 Knb


\subsection{}\label{33.3.6} Between each two principal maxima, there are N — 1
a d d ition a l m in im a twhere p may be any positive whole number with the exception
of TV, 2TV, 3TV, etc. Accordingly, there are TV—2 additional maxima
whose intensity is negligibly low compared to that of
the principal maxima.
The angular “width” of a principal maximum of the nth order, i.e. the difference between the values of the angles'  corresponding
to the additional minima bounding this maximum, is
2X   2X
Nd COS ~~ Lcostyn ’
where af>n = arc sin (nX/d), and L = Nd is the length of the
diffraction grating. For principal maxima of not particularly
high orders, angles *i|5n are small and cos   1, so that
« 2XIL.


\subsection{}\label{33.3.7} With monochromatic light the diffraction pattern on
screen E in the focal plane of lens L (Fig. 33.7) has, at large TV
values, the form of narrow and
bright principal maxima separated
by practically dark and
wide intervals. If the ratio djb
is an irrational number, the intensities
of the principal maxima
(33.3.5) decrease monotoni-
cally with an increase in the
order n of the maximum.
When the grating is illuminated
by white light, an uncoloured
central zero-order maximum is
observed on the screen. On both
sides of this maximum, diffraction
spectra of the 1st, 2nd, etc.
orders are observed. These spectra
appear as iridescent bands in
which there is a continuous transition in colour from blue-
violet at the inner edge of the spectrum to red at the outer
edge.


\subsection{}\label{33.3.8} Upon inclined incidence of light on a diffraction grating
(Fig. 33.8), the condition for the principal maxima is of the
form
d (sin — sin 0 = ±nX (n = 0, 1, 2, . . .),where i is the angle of incidence of light on the surface of the
grating. Frequently, the directions of the rays of light falling
on the grating and diffracted from it are specified by the angles
a0 and a which these rays make with the axis OX of coordinates,
arranged in the plane of the grating, perpendicular to the slits.
In this case, the condition for the principal maxima can be
rewritten in the form
d (cos a — cos a0) = ± nX.


\subsection{}\label{33.3.9} Two screens are said to be complementary if the apertures
in one of them exactly coincide in shape, size and relative position
with the opaque areas in the other screen, and vice versa.
Complementary, for example, are an opaque screen in the
form of a circle (disk) of radius R and an opaque screen having
a circular aperture of the same radius R. On the basis of the
I luygcns-Fresnel principle, it is possible to prove Babinefs
theorem (also called Babinet' s principle). This theorem states
that in Fraunhofer diffraction by any screen, the intensity of
diffracted light in any direction, except that of the propagation
of a plane wave falling on the screen, should be the same as in
diffraction by the complementary screen.


\subsection{}\label{33.3.10} Example 4. Diffraction by a large number of identical
and identically oriented obstacles. The light intensity I at an
arbitrary point M of the diffraction pattern, like in diffraction
by a one-dimensional grating (33.3.4), can be presented in the
form: I = flt. Here 7X is the intensity at point M upon diffraction
of the same incident plane wave by a single obstacle. The
function / depends only upon the number and relative arrangement
of the obstacles.
If the obstacles are arranged quite randomly with respect to one
another and their number N is large, then fzzN. In this case,
the distribution of light intensity is the same as in diffraction
by a single obstacle. But the intensity at each point of the
diffraction pattern is N times as much. For example, in diffraction
of light by a glass plate coated with spherical dust particles
of diameter <2, a system of bright concentric interference
rings is observed. The sizes of these rings correspond to the
diffraction of light by an opaque disk of diameter d or, in accordance
with Babinet’s principle (33.3.9), by a round aperture of
the same diameter,

\section{Diffraction by a Space Lattice}
\label{sec-33.4} 


\subsection{}\label{33.4.1} A space, or three-dimensional, difraction grating is an
optically nonhomogeneous medium in which the inhomogeneities
are periodically repeated in varying all three spatial coordinates.
An example of a space diffraction grating is the crystal lattice
of a crystalline solid. Particles at the lattice points (atoms,
molecules or ions) serve as orderly arranged centres that coherently
scatter light falling on them. Assume that dly d2 and d3
are the lattice constants along the three coordinate axes £, r] and £ that are oriented along three edges of the lattice intersecting
in one of its points. Then in Fraunhofer diffraction (33.2.1),
the principal maxima comply with the Laue conditions
d1 (cos a — cos a0) = n±X, d2 (cos P — cos P0) = n2X and
d3 (cos y — cos y0) = n3X.
Here a0, P0 and Yoiand a, P and y are the angles between the coordinate
axes £, r] and £, and the directions of propagation of the
incident and diffracted rays, respectively, nt, n2 and n3 are
whole numbers determining the order of the maximum, and X is
the wavelength of the light. The Laue conditions follow from the
relations (33.3.8) for diffraction maxima in inclined incidence of
light on a one-dimensional diffraction grating.

\subsection{}\label{33.4.2} Of the three angles a, P and y (and likewise for angles
a0, p0 and Yo) on y tw° angles are independent because they
must comply with one geometric relationship whose specific
form depends upon the angles between the coordinate axes £, r|
and £. For example, if the coordinate axes are mutually per­ pendicular, i.e. if the lattice is orthogonal, the geometric relationship
between a, P, and y is of the form:
cos2a + cos2P + cos2y = 1-
Therefore, with an arbitrarily given direction of incidence on
a space lattice of monochromatic light with a given wavelength X,
it is impossible, in general, to find values of a, P and y that
simultaneously comply with the geometric relationship and
Laue’s three conditions. The only exception is the zero-order
maximum (nt = n2 = n3 = 0), for which a = a0, P = p0
and Y = Yo-
In order to observe a diffraction maximum of the order (n±, n2
and n3) at given values of the angles a0, Pf and yQ, it is necessaryXU] DIFFRACTION BY A SPACE LATTICE 475
for the incident light to have a certain definite wavelength.
For instance, in the case of an orthogonal lattice, the wavelength
should be
rt\ | rt2 q . n o
c os a 0 c os P o - f - j 1- c o s y 0
X-= —2
If the wavelength X of the incident light is a fixed value, the
Laue conditions and the geometric relationship between angles a,
P and y can simultaneously be complied with by a suitable
choice of the direction of light incidence on the diffraction
grating, i.e. of angles oc0, P0 and y0.

\subsection{}\label{33.4.3} It follows from Laue’s conditions that at X   2dmax,
where dmax is the maximum of the values of d±, d2 and d3,
there should be no maxima, except the zero-order one (rci= n2 =
- rc3 = 0). Light of such wavelengths propagates in a medium
“without noticing” its inhomogeneities, i.e. without undergoing
diffraction. Hence, the condition X> 2dmax is called the
condition for optical homogeneity of the medium.
'The crystal lattice constants of solids are considerably shorter
I lian the wavelengths of visible light (dt   5 X 10 ~10 m, where-
as  vjs i   5 X 10~7 m). Hence, crystals are optically homogeneous
media for visible light.* At the same time, crystals are
natural diffraction lattices for X rays (31.4.4).

\subsection{}\label{33.4.4} The diffraction of X rays by crystals can be interpreted as
t he result of interference of X rays that are mirror reflected from
a system of parallel planes passing through the crystal lattice
points. They are called planes of atoms, or lattice planes, of the
crystal. The distance d between two adjacent lattice planes is
railed the interplanar spacing, or distance, and the angle 0 between
the incident ray and the lattice plane (Fig. 33.9) is called
I lie grazing angle. The path difference of rays reflected from two
adjacent lattice planes** is A = BC + BD = 2d sin 0. There-
* Molecular scattering of visible light (34.3.1) is possible in
crystals. It is associated with the impairment of their optical
homogeneity due to fluctuations in density.
** X rays are not refracted in crystals because the values of the
refractive index (31.5.1) of all crystals for electromagnetic radiation
of such high frequency are practically equal to unity.fore, according to Sect. 32.2.1, reflection is observed only in the
directions, corresponding to diffraction maxima that comply
with the Wulff-Bragg condition:
2d sin 0 = mk,
where m = 1, 2, . . . is the order of the diffraction maximum.

\subsection{}\label{33.4.5} The experimental method of investigating the atomic
structure of matter by analysing the diffraction of X radiation
as it passes through the substance
being investigated is
called X-ray structure analysis.
This method proves most
efficient in studying the
structure of crystalline bodies.
The diffraction pattern,
registered on photographic
film, is called the X-ray photoFig. 33.9
graph,
or pattern, of the specimen.
An X-ray photograph,
obtained in the diffraction of a beam of “white1’ X radiation
(with a continuous frequency spectrum) by a monocrystal is
called a Laue diffraction pattern. It consists of discrete diffraction
(Laue pattern) spots whose position is determined by the Laue
conditions (33.4.1). A Debye-Scherrer-IIull powder photograph,
or X-ray powder photograph, is an X-ray photograph obtained in
the diffraction of a beam of monochromatic X radiation by
a poly crystalline specimen (for instance, a crystalline powder).
Such a photographic pattern consists of a system of concentric
diffraction rings. The radii of the rings r = I tan 20, where I is
the distance from the specimen to the plane of the X-ray powder
photograph, perpendicular to the incident ray, and angles 0
comply with the Wulff-Bragg condition (33.4.4).

\section{Resolving Power of Optical Instruments}
\label{sec-33.5} 

\subsection{}\label{33.5.1} The image of the object in any optical system (telescope,
microscope, photographic camera, etc.) is obtained by means of
a restricted beam of light admitted into the instrument by the
so-called aperture diaphragm. Such a diaphragm may be that of
a photographic camera, lens mount of a telescope, etc. Reducing
the diameter of the aperture diaphragm helps to decrease vari-:a5] RESOLVING POW ER 477
ous distortions of the image resulting from the use of wide beams
of light and called geometrical aberrations of the optical system.
But, owing to diffraction of light in an optical instrument, the image of a bright point is not a point, but a bright spot surrounded
by a system of concentric interference rings (alternating
dark and bright ones in the case of monochromatic light and
iridescent in the case of white light). This phenomenon limits the
resolving power ( resolution) of the optical instrument, i.e. its
capacity to provide separate images of two points of the object
close to each other.

\subsection{}\label{33.5.2} According to Rayleigh's criterion, the images of two identical
point sources of light can still be seen separately if the
central maximum of the diffraction pattern of one falls exactly
on the first minimum of the diffraction pattern of the other. It
follows from Sect. 33.3.3 that in accordance with the Rayleigh
criterion, two close stars, observed in a telescope with monochromatic
light of the wavelength X, can be resolved (viewed
separately) if the angular separation (distance) between them is
Acp> 1.22 XID,
where D is the diameter of the objective. The quantity (Aqp)0 =
= 1.22 XID is called the angular resolution limit of the telescope.
The reciprocal quantity l/(Acp)0 is called the resolving power of
the telescope. The resolving power of a telescope increases in
proportion to the diameter of its objective. The condition for
the resolution of ’a terrestrial telescope and a photographic camera
in viewing and photographing distant items coincides with
that for an astronomical telescope.
The angular resolution limit of the eye is determined by the
diffraction of light by the pupil (D   2 mm) and the granular
structure of the retina. This limit is about one minute of arc.

\subsection{}\label{33.5.3} The resolving power of a microscope is characterized by
the quantity (AZ)0, the shortest distance between two points of
the object that are separately visible on the image. In the case
of a self-luminous object, all of whose points can be assumed to
be incoherent sources,
(AZ)0 = 0.61 0
A
where A,0 is the wavelength of the light in vacuum, A = n sin u
is the numerical aperture of the objective, n is the refractivem DIFFRACTION OF LIGHT [33
index of the medium between the object and the objective, and
u is one half of the aperture angle of the beam of light issuing
from a point of the object and entering the microscope objective
(i.e. the angle between the outermost ray entering the objective
and the optical axis). For non-self-luminous objects, the value
of (Al)0 depends upon the illumination conditions. But in this
case as well, (Al)0  X 0lA.
The resolving power of a microscope can be improved either
by reducing the wavelength X0, or by increasing the numerical
aperture A (numerical aperture is abbreviated N.A.). The former
method is employed in ultraviolet and electron microscopy,
and the latter, in the immersion microscope, in which the space
between the objective and the object is filled with a transparent liquid having a refractive index t  > 1.

\section{Holography}
\label{sec-33.6} 


\subsection{}\label{33.6.1} Holography is a technique for obtaining a three-dimensional
optical image of objects and is based on the phenomenon
of wave interference. In contrast to the conventional photographic
method, light-sensitive fine-grained photographic emulsion is
used to record the relations, not only between the amplitudes
(or their squares, i.e. the intensities) of light waves scattered by
various small portions of the object’s surface, but also between
the phases of these waves.
The essence of the holographic method is made clear in
Fig. 33.10. Photographic plate P (Fig. 33.10a) records the
interference pattern obtained in superposing wave i, scattered
by the object Q and called the signal wave, or object beam% and
wave 2, coherent with wave 1 and having fixed values of amplitude
and phase. Wave 2, called the reference wave, or reference
beam, is emitted by the same source of light that illuminates the
object and, after reflection from mirror B, falls directly on the
photographic plate P. The interference pattern recorded on the
photographic plate after it is developed is called the hologram of
object Q. It consists of a very fine and intricate pattern of alternating
interference maxima and minima in the blackening of
the photographic emulsion and, in contrast to a photographic
image, has no outward resemblance to the object.
Holograms are obtained by the interference of light at large path
differences, i.e. they require an extremely high degree of coher-ence of the light (32.2.4). For this reason lasers (39.6.8) are used
in holography as the light source.

\subsection{}\label{33.6.2} Reconstructing the image oj the object from its hologram C
(Fig. 33.106) is accomplished by passing reference wave 2,
from the same laser used to make the hologram, through it as
in the projection of a slide (diapositive). In this operation, the
orientation of the plate with the hologram with respect to the
Fig. 33.10
reference wave should also be maintained. Wave 2 is diffracted
by the hologram. As a result of this diffraction, two three-
dimensional images of the object, a virtual and a real image, are
seen. The virtual image Q' is at the same place, with respect to
the hologram, where the object Q was in taking the photograph.
This image is visible in looking through the hologram as if it
was a window. The real image Q" is on the other side of the
hologram. It seems to he in the air in front of the hologram and
is a mirror image of the object.m DIFFRACTION OF LIGHT [33
Ordinarily, the virtual image of the object is used; in visual
perception it is identical to the object itself. It not only has the property of being three-dimensional, but its perspective varies
with the position of the eye of the observer with respect to
the hologram. For example, by moving the head along the
hologram, it is possible to see behind some item located in the
foreground of the holographic image.


\subsection{}\label{33.6.3} The interference pattern at each point of the hologram is
determined by light scattered by all the points of the object.
Hence, any portion of the hologram contains information
about the whole object and can enable the image of the whole
object to be restored if, in damaging the hologram, only this
single portion remains intact. The smaller the size of the saved
portion of the damaged hologram, the less the light diffracted
by it in restoring the image. The luminance is correspondingly
reduced and the sharpness of the holographic image of the
object is impaired. Thus, a hologram has essential advantages
with respect to reliability of storage of information over a conventional
photograph or photographic negative, each element of
which contains information only about the portion of the object
depicted on it.
Holographic recording of information is distinguished for its
high capacity and compact form. A great many holograms can
be recorded on a single photographic plate. For instance, it is
sufficient, for this purpose, to record each hologram with
a different value of the angle of incidence of the reference wave
on the photographic plate.

\subsection{}\label{33.6.4} Holography can be used to obtain coloured three-dimensional
images of objects. To make such a “coloured” hologram,
use is made of the monochromatic light of three lasers with
three primary colours (for instance, red, green and blue).
The recording of the interference patterns, corresponding to
light of three wavelengths, is accomplished simultaneously or
consecutively on the same photographic plate. To reconstruct
the coloured three-dimensional image of the object it is necessary
to simultaneously direct onto the hologram the three reference
beams of monochromatic light that were used in recording.

\subsection{}\label{33.6.5} Three-dimensional holograms have special properties not
possessed by the ordinary kind. They are obtained by using
thick-layer photographic emulsions. Recorded on a three-
dimensional hologram is a spatial interference pattern, instead
of a plane one. This pattern is obtained in superposing the sig-INTERACTION OR LIGHT 481
nal and reference waves. Such a hologram is similar to a spatial
diffraction grating (lattice). It is capable of separating out of the
incident white light, monochromatic light of the wavelength
that was used in recording the hologram. Therefore, an image
recorded as a three-dimensional hologram can be reconstructed
by illuminating the hologram either by the corresponding monochromatic
light or by white light. If a three-dimensional hologram
is “coloured”, the coloured three-dimensional image can
be reconstructed by simply illuminating the hologram with
white light.